\documentclass[12pt,CJK]{article}
\usepackage{geometry}
\usepackage{CJKutf8}
\usepackage{amsmath}
\usepackage{amssymb}
\usepackage{graphicx}
\usepackage{eufrak}
\usepackage{color}
\usepackage{slashed}
\usepackage{simplewick}
\usepackage{tikz}
\usepackage{tcolorbox}
\graphicspath{{../figures/}}
%%figures
\def\lfig#1#2{\includegraphics[width=#1 in]{#2}}
\def\addfig#1#2{\begin{center}\includegraphics[width=#1 in]{#2}\end{center}}
\def\wulian{\includegraphics[width=0.18in]{emoji_wulian.jpg}}
\def\bigwulian{\includegraphics[width=0.35in]{emoji_wulian.jpg}}
\def\bye{\includegraphics[width=0.18in]{emoji_bye.jpg}}
\def\bigbye{\includegraphics[width=0.35in]{emoji_bye.jpg}}
\def\huaixiao{\includegraphics[width=0.18in]{emoji_huaixiao.jpg}}
\def\bighuaixiao{\includegraphics[width=0.35in]{emoji_huaixiao.jpg}}
\def\jianxiao{\includegraphics[width=0.18in]{emoji_jianxiao.jpg}}
\def\bigjianxiao{\includegraphics[width=0.35in]{emoji_jianxiao.jpg}}
%% colors
\def\blacktext#1{{\color{black}#1}}
\def\bluetext#1{{\color{blue}#1}}
\def\redtext#1{{\color{red}#1}}
\def\darkbluetext#1{{\color[rgb]{0,0.2,0.6}#1}}
\def\skybluetext#1{{\color[rgb]{0.2,0.7,1.}#1}}
\def\cyantext#1{{\color[rgb]{0.,0.5,0.5}#1}}
\def\greentext#1{{\color[rgb]{0,0.7,0.1}#1}}
\def\darkgray{\color[rgb]{0.2,0.2,0.2}}
\def\lightgray{\color[rgb]{0.6,0.6,0.6}}
\def\gray{\color[rgb]{0.4,0.4,0.4}}
\def\blue{\color{blue}}
\def\red{\color{red}}
\def\green{\color{green}}
\def\darkgreen{\color[rgb]{0,0.4,0.1}}
\def\darkblue{\color[rgb]{0,0.2,0.6}}
\def\skyblue{\color[rgb]{0.2,0.7,1.}}
%%control
\def\be{\begin{equation}}
\def\ee{\nonumber\end{equation}}
\def\bea{\begin{eqnarray}}
\def\eea{\nonumber\end{eqnarray}}
\def\bch{\begin{CJK}{UTF8}{gbsn}}
\def\ech{\end{CJK}}
\def\bitem{\begin{itemize}}
\def\eitem{\end{itemize}}
\def\bcenter{\begin{center}}
\def\ecenter{\end{center}}
\def\bex{\begin{minipage}{0.2\textwidth}\includegraphics[width=0.6in]{jugelizi.png}\end{minipage}\begin{minipage}{0.76\textwidth}}
\def\eex{\end{minipage}}
\def\chtitle#1{\frametitle{\bch#1\ech}}
\def\bmat#1{\left(\begin{array}{#1}}
\def\emat{\end{array}\right)}
\def\bcase#1{\left\{\begin{array}{#1}}
\def\ecase{\end{array}\right.}
\def\bmini#1{\begin{minipage}{#1\textwidth}}
\def\emini{\end{minipage}}
\def\tbox#1{\begin{tcolorbox}#1\end{tcolorbox}}
\def\pfrac#1#2#3{\left(\frac{\partial #1}{\partial #2}\right)_{#3}}
%%symbols
\def\bropt{\,(\ \ \ )}
\def\optlist#1#2#3{(A)\;#1\ \ \ (B)\;#2\ \ \ (C)\;#3}
\def\sone{$\star$}
\def\stwo{$\star\star$}
\def\sthree{$\star\star\star$}
\def\sfour{$\star\star\star\star$}
\def\sfive{$\star\star\star\star\star$}
\def\rint{{\int_\leftrightarrow}}
\def\roint{{\oint_\leftrightarrow}}
\def\stdHf{{\textit{\r H}_f}}
\def\deltaH{{\Delta \textit{\r H}}}
\def\ii{{\dot{\imath}}}
\def\skipline{{\vskip0.1in}}
\def\skiplines{{\vskip0.2in}}
\def\lagr{{\mathcal{L}}}
\def\hamil{{\mathcal{H}}}
\def\vecv{{\mathbf{v}}}
\def\vecx{{\mathbf{x}}}
\def\vecy{{\mathbf{y}}}
\def\veck{{\mathbf{k}}}
\def\vecp{{\mathbf{p}}}
\def\vecn{{\mathbf{n}}}
\def\vecA{{\mathbf{A}}}
\def\vecP{{\mathbf{P}}}
\def\vecsigma{{\mathbf{\sigma}}}
\def\hatJn{{\hat{J_\vecn}}}
\def\hatJx{{\hat{J_x}}}
\def\hatJy{{\hat{J_y}}}
\def\hatJz{{\hat{J_z}}}
\def\hatj#1{\hat{J_{#1}}}
\def\hatphi{{\hat{\phi}}}
\def\hatq{{\hat{q}}}
\def\hatpi{{\hat{\pi}}}
\def\vel{\upsilon}
\def\Dint{{\mathcal{D}}}
\def\adag{{\hat{a}^\dagger}}
\def\bdag{{\hat{b}^\dagger}}
\def\cdag{{\hat{c}^\dagger}}
\def\ddag{{\hat{d}^\dagger}}
\def\hata{{\hat{a}}}
\def\hatb{{\hat{b}}}
\def\hatc{{\hat{c}}}
\def\hatd{{\hat{d}}}
\def\hatN{{\hat{N}}}
\def\hatH{{\hat{H}}}
\def\hatp{{\hat{p}}}
\def\Fup{{F^{\mu\nu}}}
\def\Fdown{{F_{\mu\nu}}}
\def\newl{\nonumber \\}
\def\vece{\mathrm{e}}
\def\calM{{\mathcal{M}}}
\def\calT{{\mathcal{T}}}
\def\calR{{\mathcal{R}}}
\def\barpsi{\bar{\psi}}
\def\baru{\bar{u}}
\def\barv{\bar{\upsilon}}
\def\qeq{\stackrel{?}{=}}
\def\torder#1{\mathcal{T}\left(#1\right)}
\def\rorder#1{\mathcal{R}\left(#1\right)}
\def\contr#1#2{\contraction{}{#1}{}{#2}#1#2}
\def\trof#1{\mathrm{Tr}\left(#1\right)}
\def\trace{\mathrm{Tr}}
\def\comm#1{\ \ \ \left(\mathrm{used}\ #1\right)}
\def\tcomm#1{\ \ \ (\text{#1})}
\def\slp{\slashed{p}}
\def\slk{\slashed{k}}
\def\calp{{\mathfrak{p}}}
\def\veccalp{\mathbf{\mathfrak{p}}}
\def\Tthree{T_{\tiny \textcircled{3}}}
\def\pthree{p_{\tiny \textcircled{3}}}
\def\dbar{{\,\mathchar'26\mkern-12mu d}}
\def\erf{\mathrm{erf}}
\def\const{\mathrm{constant}}
\def\pheat{\pfrac p{\ln T}V}
\def\vheat{\pfrac V{\ln T}p}
%%units
\def\fdeg{{^\circ \mathrm{F}}}
\def\cdeg{^\circ \mathrm{C}}
\def\atm{\,\mathrm{atm}}
\def\angstrom{\,\text{\AA}}
\def\SIL{\,\mathrm{L}}
\def\SIkm{\,\mathrm{km}}
\def\SIyr{\,\mathrm{yr}}
\def\SIGyr{\,\mathrm{Gyr}}
\def\SIV{\,\mathrm{V}}
\def\SImV{\,\mathrm{mV}}
\def\SIeV{\,\mathrm{eV}}
\def\SIkeV{\,\mathrm{keV}}
\def\SIMeV{\,\mathrm{MeV}}
\def\SIGeV{\,\mathrm{GeV}}
\def\SIcal{\,\mathrm{cal}}
\def\SIkcal{\,\mathrm{kcal}}
\def\SImol{\,\mathrm{mol}}
\def\SIN{\,\mathrm{N}}
\def\SIHz{\,\mathrm{Hz}}
\def\SIm{\,\mathrm{m}}
\def\SIcm{\,\mathrm{cm}}
\def\SIfm{\,\mathrm{fm}}
\def\SImm{\,\mathrm{mm}}
\def\SInm{\,\mathrm{nm}}
\def\SImum{\,\mathrm{\mu m}}
\def\SIJ{\,\mathrm{J}}
\def\SIW{\,\mathrm{W}}
\def\SIkJ{\,\mathrm{kJ}}
\def\SIs{\,\mathrm{s}}
\def\SIkg{\,\mathrm{kg}}
\def\SIg{\,\mathrm{g}}
\def\SIK{\,\mathrm{K}}
\def\SImmHg{\,\mathrm{mmHg}}
\def\SIPa{\,\mathrm{Pa}}

\geometry{tmargin=0.3in, bmargin=0.5in, lmargin=0.5in, rmargin=0.9in, nohead, nofoot}
\def\mark#1{{\color{blue} (#1分)}}
\renewcommand{\thepage}{}
\begin{document}
\bch
{\large 数理方法 课堂小测I}

{\vskip 0.2in}

姓名 ....................... {\hskip 0.5in}    学号 .......................{\hskip 0.5in}  分数 ...................

\bitem
\item[(一)]{选择题，每题3分，共45分。

  \bitem
\item[(1)]{ $e^{\ii \pi}$的值为  \bropt
  
  \optlist{$-1$}{$0$}{$1$}{$\ii$} }
\item[(2)]{方程$e^z = 1 +\ii$ 的全部复数解为 $z=$ \bropt


\foptlist{$\ln 2 + \frac{\pi}{4}\ii$}{$\ln 2 + \left(2n+\frac{1}{4}\right)\pi \ii,\,n\in Z$}{$\frac{1}{2}\ln 2 + \left(2n+\frac{1}{4}\right)\pi \ii,\,n\in Z$}{$\frac{1}{2}\ln 2 + \left(n+\frac{1}{4}\right)\pi \ii,\,n\in Z$}}


\item[(3)]{复变函数 $z\cos z$ 在 $z=0$ 处的导数为 \bropt
  
  \optlist{$-1$}{$0$}{$1$}{$2\pi\ii$}}

\item[(4)]{复变函数$f(z) = (1+z)e^z$的原函数为（忽略不写积分常数）： \bropt

  \optlist{$e^z$}{$ze^z$}{$\left(z+\frac{z^2}{2}\right)e^z$}{$\frac{e^z}{1+z}$}}
\item[(5)]{方程 $z^5+ z^4+z^3+2=0$ 的所有复数解的平方和为 \bropt

  \optlist{$-1$}{$0$}{$1$}{$2$}}  

\item[(6)]{已知$f$和$g$在全复平面上解析，且在某个复数集$S$上$f$和$g$恒等。当$S$为下述哪个集合时，我们{\bf 不能}断言$f$和$g$在全复平面上恒等? \bropt

  \optlist{单位圆$|z|=1$内部}{区间$(0,1)$}{$(0,1)$上的所有有理数}{整数集}}
  
\item[(7)]{$\frac{1}{(1+e^z)\sin z}$ 在区域$|z|< 5 $内有多少个孤立奇点? \bropt
  
  \optlist{$3$}{$4$}{$5$}{$6$}}
\item[(8)]{$\frac{1}{z^2-3z+2}$ 在 $z=1$ 处的留数等于 \bropt

    \optlist{$2$}{$1$}{$0$}{$-1$}}  

\item[(9)]{下列哪个多值函数在区域 $1<|z|<2$ 内可以规定适当的幅角范围成为解析函数？ \bropt

  \optlist{$\ln(z-1)$}{$\ln (z+1)$}{$\ln [(z-1)(z+1)]$}{$\ln \frac{z-1}{z+1}$}}


\item[(10)]{$\sqrt{2\pi}\,\delta (x^2-1)$的傅立叶变换为 \bropt

  \optlist{$1$}{$\cos k$}{$\sin k$}{$\frac{\sin k}{k}$}}

\item[(11)]{$e^{-t}\sin{2t}$的拉普拉斯变换为 \bropt

  \optlist{$\frac{2}{(p+1)^2+4}$}{$\frac{2}{(p-1)^2+4}$}{$\frac{2p}{p^2+4}$}{$\frac{2e^{-p}}{p^2+4}$}}

\item[(12)]{$ \left\vert\oint_{|z|=1}\, \frac{\cos z}{z^{2719}}\,dz\right\vert$ 和下列哪个数量级最接近？ \bropt

  \optlist{$10^{-900}$}{$10^{-2700}$}{$10^{-5100}$}{$10^{-15300}$}}

\item[(13)]{ 在单位球内的积分 $\iiint_{x^2+y^2+z^2\le 1} (|x|+|y|+|z|)dxdydz$ 等于\bropt

  \optlist{$\frac{3}{2}\pi$}{$\sqrt{3}\pi$}{$2\pi$}{$2\sqrt{3}\pi$}}
  
\item[(14)]{用$z^*$表示$z$的共轭复数，按逆时针沿着曲线$|z-3|+|z+3|=10$的围道积分$\oint_{|z-3|+|z+3|=10}\, (z^*dz - zdz^*)$  等于 \bropt
  
  \optlist{$48\pi\ii$}{$60\pi\ii$}{$80\pi\ii$}{$96\pi\ii$}}

\item[(15)]{下列哪一个复变函数在复平面上处处不可导？ \bropt

  \optlist{$\frac{1}{z}$}{$|z|^2$}{$\sin{|z|}$}{$e^{|z|}$} }  
    
  \eitem
  
\vspace{1in}
  }
\item[(二)]{ 填空题，每题5分，共35分。

  \bitem  
\item[(1)]{复变函数 $e^{\frac{\cos z}{1+z^2+z^4}}\cos{(z^3)}$ 在 $z=0$ 处的导数为\_\_\_\_\_\_\_\_。 }  
\item[(2)]{函数 $\frac{1}{z^{5}+z+1}$ 的所有孤立奇点处的留数之和为\_\_\_\_\_\_\_\_。}  
\item[(3)]{积分 $\int_0^{\frac{\pi}{2}}\, \cos^5\theta \, \sin^4\theta\, d\theta$ 等于\_\_\_\_\_\_\_\_。   }
\item[(4)]{实积分 $\int_0^\infty e^{-x^2}\cos{2x}dx$ 等于\_\_\_\_\_\_\_\_。}
\item[(5)]{逆时针方向沿着上半个单位圆的积分$\int_{|z|=1, \mathrm{Im}(z)\ge 0}\, z^{\frac{1}{3}}\,dz$ 等于\_\_\_\_\_\_\_\_。}
\item[(6)]{复变函数$f(z) = |1-z|^4 + (1+|z|)^4$的最小可能值等于\_\_\_\_\_\_\_\_。}   
\item[(7)]{定义函数  $f(x) = \frac{1}{\pi}\int_0^\pi \cos{\left( x \cos t\right)} \, dt. $，则积分 $\int_0^{\frac{\pi}{2}} f(x)f(\frac{\pi}{2}-x)\,dx$ 等于 \_\_\_\_\_\_\_\_。}
\eitem

  }
  

\item[(三)]{长度为$L$的不良导体棒一端和温度为$T_0$的热库接触，并在$t=0$时刻和热库处于热平衡。从$t=0$时刻开始，在不良导体棒的另一端注入恒定大小为$j$的热流。设不良导体棒的导热系数$\lambda$，单位质量的比热$c$和质量密度$\rho$均已知。写出不良导体棒上温度$T(x, t)$ ($0\le x\le L, t\ge 0$)满足的方程和边界条件。(10分) 并简要分析当$t$很大时的解的渐近行为。(10分)

}
\eitem    


\ech
\end{document}
